% Mechanical relocation generated from the preservation ledger.
% Existing prose is included byte-for-byte from preserved blocks.

\subsection{Entropy Models}

\subsubsection{Model Card: Entropic Transport}

\paragraph{Formulation and required information.}
Balanced entropic transport allocates flow over feasible OD support while
matching supplied journey-origin and journey-destination marginals and trading
generalized cost against entropy.  The unbalanced variant replaces exact
marginal equality with declared divergence penalties.  Both require a feasible
support, costs, a regularization scale, and genuine passenger-journey
marginals---not raw boarding and alighting totals in a network with transfers.

\paragraph{Assumptions and limitations.}
The resulting table is selected by the chosen cost, entropy scale, support, and
marginal penalties; it is not identified by APC counts alone.  Balanced
transport is infeasible when supported marginals are incompatible.  Unbalanced
transport can reconcile disagreement, but the reconciliation is a modeling
choice and does not change leg-event counts into journey marginals.

\paragraph{Appropriate use.}
Entropy models are valuable when credible journey marginals exist, for
synthetic validation, or as a smooth initializer for another demand model.
They are auxiliary baselines rather than the primary estimator in the present
full-network setting.

For feasible support $C$, generalized cost $c_{od}$, and entropy scale
$\varepsilon>0$, the balanced model solves
\[
 \min_{x\geq0}
 \sum_{(o,d)\in C}c_{od}x_{od}
 +\varepsilon\sum_{(o,d)\in C}x_{od}(\log x_{od}-1)
\]
subject to supplied journey-origin and journey-destination marginals. The
unbalanced variant replaces exact marginal constraints with declared positive
Kullback--Leibler penalties. It can reconcile inconsistent totals, but cannot
turn vehicle-leg event counts into passenger-journey marginals.
